\section{Introduction}

Online streaming platforms depend on recommender systems to surface relevant items from catalogs that users
will never fully browse.
The problem is difficult because user--item interaction matrices are extremely sparse: most users rate only
a small fraction of available movies, and popularity follows a long-tail distribution where niche titles
receive few ratings.

\textbf{MovieLens 1M} provides approximately one million explicit star ratings from roughly 6{,}000 users
on about 3{,}900 movies, along with user demographics and movie metadata \cite{harper2015movielens}.
It is a standard benchmark for rating prediction and hybrid recommendation research.

This capstone addresses three practical goals:
\begin{enumerate}
  \item Understand data characteristics that affect model choice (sparsity, demographics, cold start).
  \item Build and compare baseline, machine learning, and deep learning predictors under a strict
        time-based evaluation protocol.
  \item Deliver reproducible artifacts: a single Jupyter notebook for training and evaluation, plus
        evidence tables committed to the repository.
\end{enumerate}

Accurate offline evaluation matters because deploying a weak model wastes user attention, while
overfitting to random splits inflates reported metrics.
We therefore split data chronologically by rating timestamp so that training never sees future
interactions---a setting closer to production deployment than a random shuffle.
